\documentclass[conference]{IEEEtran}

\usepackage{cite}
\usepackage{amsmath,amssymb,amsfonts}
\usepackage{graphicx}
\usepackage{booktabs}
\usepackage{array}
\usepackage{url}
\usepackage{xcolor}
\usepackage{algorithm}
\usepackage{algpseudocode}

\title{Notes on Neurons and How They Train a Model}

\author{
\IEEEauthorblockN{Luis Mendez}
}

\begin{document}
\maketitle

\textbf{Question:} How does a neuron work, and how do we use neurons to
train a model?

\bigskip

\textbf{Answer:}

\bigskip

\textbf{1. What a Neuron Computes}

An artificial neuron is a loose model of a biological neuron: dendrites
receive signals, the cell body combines them, and the axon fires an
output if the combined signal is strong enough. The artificial version
does the same thing with numbers.

Given an input vector $\mathbf{x} = (x_1, \dots, x_m)$, a neuron holds one
\textbf{weight} $w_i$ per input and one \textbf{bias} $b$. It first forms
a weighted sum,

\[
  z = \sum_{i=1}^{m} w_i x_i + b = \mathbf{w}^\top \mathbf{x} + b,
\]

then passes $z$ through a nonlinear \textbf{activation function} $\phi$ to
produce its output,

\[
  a = \phi(z).
\]

\begin{itemize}
    \item $w_i$: how strongly input $x_i$ influences the neuron. Learned
          during training.
    \item $b$: shifts $z$ up or down independently of the inputs (an
          always-on input, sometimes written as $w_0 x_0$ with $x_0=1$).
          Also learned.
    \item $\phi$: fixed in advance (not learned); it decides \emph{how}
          the weighted sum turns into an output.
\end{itemize}

\bigskip

\textbf{2. Why the Activation Function Matters}

If $\phi$ were the identity, stacking neurons into layers would still
only compute a linear function of $\mathbf{x}$ --- a composition of
linear maps is itself linear, no matter how many layers are stacked.
Nonlinear $\phi$ is what lets a network represent curved decision
boundaries and complex relationships.

\begin{itemize}
    \item \textbf{Sigmoid:} $\phi(z) = \dfrac{1}{1+e^{-z}} \in (0,1)$.
          Reads as a probability; common at the output layer for binary
          classification.
    \item \textbf{ReLU:} $\phi(z) = \max(0, z)$. Cheap, gradient is $0$
          or $1$, the default for hidden layers.
    \item \textbf{tanh:} $\phi(z) = \dfrac{e^z-e^{-z}}{e^z+e^{-z}} \in
          (-1,1)$. Like sigmoid but zero-centered.
    \item \textbf{Softmax:} generalizes sigmoid to multiple classes,
          producing a probability distribution over $K$ outputs.
\end{itemize}

\bigskip

\textbf{3. From One Neuron to a Network}

A single neuron can only separate data with a straight line (or
hyperplane). Layering neurons and connecting every neuron in one layer
to every neuron in the next (a fully-connected, or \emph{dense}, layer)
lets the network build up complex features. For layer $\ell$,

\[
  \mathbf{z}^{(\ell)} = W^{(\ell)} \mathbf{a}^{(\ell-1)} + \mathbf{b}^{(\ell)},
  \qquad
  \mathbf{a}^{(\ell)} = \phi\!\left(\mathbf{z}^{(\ell)}\right),
\]

with $\mathbf{a}^{(0)} = \mathbf{x}$. Running this input-to-output is
\textbf{forward propagation}: it turns raw features into a prediction
$\hat{y} = \mathbf{a}^{(L)}$. Each hidden layer re-represents the data;
deeper layers combine simple patterns from earlier layers into more
abstract ones.

\bigskip

\textbf{4. Measuring How Wrong the Network Is}

Training means choosing every $W^{(\ell)}, \mathbf{b}^{(\ell)}$ so that
predictions $\hat{y}$ are close to true targets $y$. A \textbf{loss
function} $C$ quantifies the gap, e.g.

\[
  C = \frac{1}{2}(\hat{y}-y)^2 \quad \text{(regression)}, \qquad
  C = -\big(y\log p + (1-y)\log(1-p)\big) \quad \text{(binary classification)},
\]

where $p$ is the sigmoid output. Training is the optimization problem of
minimizing $C$, averaged over the training set, as a function of every
weight and bias.

\bigskip

\textbf{5. Gradient Descent: How Weights Actually Move}

$C$ is a nonconvex function of thousands (or millions) of parameters, so
there is no closed-form minimizer. \textbf{Gradient descent} instead
nudges every parameter opposite its gradient, the direction of steepest
increase, so that $C$ decreases:

\[
  \mathbf{w} \gets \mathbf{w} - \eta \, \nabla_{\mathbf{w}} C,
\]

where $\eta$ is the \textbf{learning rate}. Too large an $\eta$ overshoots
and can diverge; too small makes training crawl and risks stalling in a
shallow local minimum.

\bigskip

\textbf{6. Backpropagation: Computing the Gradient Efficiently}

Backpropagation computes $\nabla_{\mathbf{w}} C$ for \emph{every} weight
in one backward sweep, by reusing the chain rule layer by layer instead
of recomputing each partial derivative from scratch. Define the error at
layer $\ell$ as $\boldsymbol{\delta}^{(\ell)} = \partial C /
\partial \mathbf{z}^{(\ell)}$. At the output layer,

\[
  \boldsymbol{\delta}^{(L)} = \frac{\partial C}{\partial \mathbf{a}^{(L)}} \odot \phi'\!\left(\mathbf{z}^{(L)}\right),
\]

and each earlier layer's error is derived from the layer after it,

\[
  \boldsymbol{\delta}^{(\ell)} = \left(\left(W^{(\ell+1)}\right)^{\top} \boldsymbol{\delta}^{(\ell+1)}\right) \odot \phi'\!\left(\mathbf{z}^{(\ell)}\right).
\]

Once every $\boldsymbol{\delta}^{(\ell)}$ is known, the weight and bias
gradients fall out directly:

\[
  \frac{\partial C}{\partial W^{(\ell)}} = \boldsymbol{\delta}^{(\ell)} \left(\mathbf{a}^{(\ell-1)}\right)^{\top},
  \qquad
  \frac{\partial C}{\partial \mathbf{b}^{(\ell)}} = \boldsymbol{\delta}^{(\ell)}.
\]

\begin{algorithm}[h]
\caption{Training one example}
\begin{algorithmic}[1]
\State \textbf{Forward pass:} compute $\mathbf{z}^{(\ell)}, \mathbf{a}^{(\ell)}$ for $\ell=1,\dots,L$
\State \textbf{Output error:} $\boldsymbol{\delta}^{(L)} \gets (\hat{y}-y)\odot\phi'(\mathbf{z}^{(L)})$
\For{$\ell = L-1,\dots,1$}
  \State $\boldsymbol{\delta}^{(\ell)} \gets \left((W^{(\ell+1)})^\top \boldsymbol{\delta}^{(\ell+1)}\right)\odot\phi'(\mathbf{z}^{(\ell)})$
\EndFor
\For{$\ell = 1,\dots,L$}
  \State $W^{(\ell)} \gets W^{(\ell)} - \eta\,\boldsymbol{\delta}^{(\ell)}(\mathbf{a}^{(\ell-1)})^\top$
  \State $\mathbf{b}^{(\ell)} \gets \mathbf{b}^{(\ell)} - \eta\,\boldsymbol{\delta}^{(\ell)}$
\EndFor
\end{algorithmic}
\end{algorithm}

\bigskip

\textbf{7. Batch vs. Stochastic vs. Mini-Batch}

In practice, the gradient is not computed from a single example or the
whole dataset at once:

\begin{itemize}
    \item \textbf{Batch GD}: exact gradient over all $n$ examples per
          update; accurate but slow, and deterministic enough to get
          stuck in a local minimum.
    \item \textbf{Stochastic GD (SGD)}: one example per update; noisy but
          cheap, and the noise helps escape local minima/saddle points.
    \item \textbf{Mini-batch GD}: a small batch (e.g. 32--256 examples)
          per update; the standard in practice, balancing gradient
          variance against compute cost and vectorizing well on a GPU.
\end{itemize}

\bigskip

\textbf{8. Putting It Together}

\begin{itemize}
    \item A neuron = weighted sum + bias, squashed by an activation
          function.
    \item Stacking neurons into layers (forward propagation) turns raw
          inputs into a prediction.
    \item A loss function scores how wrong that prediction is.
    \item Backpropagation computes how each weight should change to
          reduce the loss.
    \item Gradient descent (batch, stochastic, or mini-batch) actually
          applies that change, repeated over many epochs, until the
          model fits the training data.
\end{itemize}

\bigskip

\textbf{9. Typical Keras Example}

\begin{verbatim}
model = Sequential([
    Dense(64, activation='relu', input_shape=(m,)),
    Dense(32, activation='relu'),
    Dense(1, activation='sigmoid')
])

model.compile(
    optimizer='adam',
    loss='binary_crossentropy',
    metrics=['accuracy']
)

model.fit(X_train, y_train, epochs=50, batch_size=32)
\end{verbatim}

Each \texttt{Dense} layer is a bank of neurons; \texttt{fit} runs forward
propagation, computes the loss, backpropagates the gradient, and applies
mini-batch updates (via \texttt{adam}, an adaptive variant of gradient
descent) once per batch, for the given number of epochs.

\end{document}
